\begin{textsample}{POS Dim 3 – vem\_ed\_14 – Score 5.00 – vem\_ed\_14/t058.txt}  \label{ex:f3_pos_088}
EDITORIAL Em respeito à lei eleitoral, houve uma pausa nas publicações do Centro Paula Souza entre julho e \textbf{outubro} de 2022.
Por isso, esta edição circula somente após o fim do pleito estadual.
Durante o período sem publicações externas, manteve-se a comunicação interna entre a equipe dos PCIs / Cesu e os professores e equipes \textbf{gestoras} das Fatecs.
Além das postagens na equipe Conexão PCI / Cesu na Plataforma Teams, foi criado um novo canal de comunicação interna em agosto de 2022:o grupo de WhatsApp PCI Informações. \textbf{Caso} tenha interesse em participar do grupo, segue o link: https://bit.ly/3sYiVx.
Uma notícia que traz muito orgulho: em julho de 2022, as Fatecs do Centro Paula Souza se destacam entre as instituições com maior número e projetos COIL (Intercâmbios Virtuais) no mundo, segundo o site COIL Connect for Virtual Exchange (https://coilconnect.org/).
Apenas no último ano (agosto 2021–julho 2022, seguindo como base o ano letivo americano), foram 109 PCIs: 39 no segundo semestre de 2021 e 70 projetos no primeiro semestre de 2022.
A percepção de alunos e professores que realizam os PCIs é muito satisfatória, como se vê nas páginas 4 e 5.
Os PCIs também têm gerado publicações acadêmicas e participações em congressos internacionais.
Confira também o depoimento de Adya Sharma, diretora do Symbiosis Centre for Management Studies (Pune, Índia).
Boa leitura!
EXPEDIENTE Expediente CPS Diretora-Superintendente: Laura Laganá Vice-Diretora-Superintendente: Emilena Lorenzon Bianco Chefe de Gabinete: Armando Natal Maurício Expediente Cesu Coordenador Técnico: Rafael Ferreira Alves Diretor Acadêmico-Pedagógico: André Luiz Braun Galvão Departamento Administrativo: Elisete Buttignon EDI - Estruturação e Desenvolvimento Instrucional: Thais Lari Braga Cilli Expediente Línguas e Projetos Colaborativos Internacionais - Cesu Coordenação de Línguas e Projetos Internacionais: Mariane Teixeira Coordenação de Projetos Colaborativos Internacionais: Osvaldo Succi Junior Acompanhamento \textbf{pedagógico} PCI: Ana Carolina Freschi, Neusa Haruka Gritti e Regiane Moreira Expediente VEm Corpo editorial: Ana Carolina Freschi, Mariane Teixeira, Neusa Haruka Gritti, Osvaldo Succi Junior e Regiane Moreira Jornalista \textbf{responsável} e Comunicação: Patrícia Patrício - MTb 25.131 Editoração e diagramação: Fábio Gomes da Silva VEm: Virtual Exchange Medium é um informativo com publicação bimestral da Cesu / CEETEPS: Rua dos Andradas, 140 - Santa Efigênia - 01208-000 - São Paulo - SP

% matched lemmas: caso, gestora, outubro, pedagógico, responsável
\end{textsample}
